\chapter{Giới thiệu}

\section{Bối cảnh và mục tiêu}
Đề tài của môn Lập trình Python là \emph{Ứng dụng dự đoán giá bằng Machine Learning}. Báo cáo giữa kỳ thực hiện hai phần đầu của đề tài trên dữ liệu giao dịch nhà ở Melbourne (Úc): \emph{phân tích dữ liệu} (nguồn, kiểm tra, làm sạch, thống kê, trực quan hóa) và \emph{mô hình hồi quy} dự đoán giá bằng scikit-learn. Dữ liệu gồm 13,580 căn nhà bán từ 28/01/2016 đến 23/09/2017 (\bs{N01, N02}). Giao diện người dùng, cơ sở dữ liệu và phân quyền thuộc đồ án riêng, không nằm trong báo cáo này.

Các mục tiêu cụ thể:
\begin{enumerate}[nosep]
  \item Mô tả nguồn gốc, quy mô và ý nghĩa các cột của dữ liệu (Chương~\ref{ch:nguon}).
  \item Kiểm tra chất lượng dữ liệu và làm sạch bằng các quy tắc có lý do, có số dòng bị ảnh hưởng (Chương~\ref{ch:lamsach}).
  \item Thực hiện bảy phép phân tích thống kê gồm mục đích, phương pháp, kết quả và diễn giải (Chương~\ref{ch:phantich}).
  \item Trình bày năm biểu đồ thuộc năm dạng khác nhau (Chương~\ref{ch:trucquan}).
  \item Trình bày cơ sở lý thuyết của các thành phần scikit-learn được dùng: \texttt{Pipeline}, bộ tiền xử lý, Ridge, Random Forest, chia tập và chỉ số đánh giá (Chương~\ref{ch:lythuyet}).
  \item Huấn luyện, chọn và đánh giá mô hình hồi quy so với một mô hình cơ sở (baseline) (Chương~\ref{ch:mohinh}).
  \item Rút ra nhận xét có số liệu, nêu hạn chế và hướng cải thiện (Chương~\ref{ch:ketluan}).
\end{enumerate}

\section{Phạm vi và câu hỏi nghiên cứu}
\paragraph{Phạm vi.} Báo cáo chỉ dùng tệp \texttt{data/melb\_data.csv}. Mã nguồn gồm \texttt{pipeline.py} (kiểm tra, làm sạch, phân tích, biểu đồ), \texttt{train\_model.py} (huấn luyện và đánh giá mô hình) và notebook \texttt{melbourne\_\allowbreak data\_\allowbreak analysis\_\allowbreak pipeline.ipynb} (trình bày từng bước). Kết luận chỉ áp dụng cho Melbourne giai đoạn 2016--2017.

\paragraph{Nguồn số liệu.} Mọi số liệu trong báo cáo thuộc một trong ba loại. Loại thứ nhất là các tệp do \texttt{pipeline.py} và \texttt{train\_model.py} ghi vào \texttt{outputs/}. Loại thứ hai là các số tính thêm, do hàm \texttt{build\_report\_metrics()} (\texttt{pipeline.py}) và \texttt{build\_model\_error\_metrics()} (\texttt{train\_model.py}) ghi vào \texttt{outputs/report\_metrics.csv} và \texttt{outputs/model\_error\_metrics.csv} cùng biểu thức tính. Loại thứ ba là các kết quả đối chứng trong \texttt{outputs/doi\_chung/}. Dưới mỗi bảng và hình có dòng \emph{Nguồn} chỉ rõ tệp hoặc mã số liệu; Phụ lục~\ref{pl:nguon} tổng hợp cách tái lập.

\paragraph{Câu hỏi nghiên cứu.}
\begin{description}[leftmargin=2.2em,style=sameline,nosep]
  \item[Q1] Giá nhà phân phối thế nào, có lệch hoặc có ngoại lai đáng kể không? (Phép 1)
  \item[Q2] Giá khác nhau ra sao giữa các loại nhà, vùng và suburb? (Phép 2, 3, 5)
  \item[Q3] Số phòng và khoảng cách đến trung tâm ảnh hưởng đến giá thế nào? (Phép 4, 6)
  \item[Q4] Thuộc tính số nào tương quan mạnh nhất với giá, và ngoại lai làm sai lệch điều đó ra sao? (Phép 6)
  \item[Q5] Mô hình dự đoán tốt đến đâu so với baseline và yếu ở phân khúc nào? (Chương~\ref{ch:mohinh})
  \item[Q6] Vì sao mỗi thành phần scikit-learn được chọn và hoạt động thế nào? (Chương~\ref{ch:lythuyet})
  \item[Q7] Dữ liệu và cách làm có những hạn chế nào? (mục~\ref{sec:hanche})
\end{description}

\section{Công cụ và thư viện}
Toàn bộ xử lý viết bằng Python; các thư viện chính lấy từ \texttt{requirements.txt} (Bảng~\ref{tab:thuvien}).

\begin{table}[H]
\centering
\caption{Thư viện chính và vai trò trong dự án}\label{tab:thuvien}
\small
\begin{tabular}{@{}l L{6.2cm} L{4.6cm}@{}}
\toprule
Thư viện & Vai trò & Nơi dùng \\
\midrule
NumPy & Tính toán số, biến đổi \texttt{log1p} và ngược lại cho giá & \texttt{pipeline.py}, \texttt{train\_model.py} \\
Pandas & Đọc CSV, kiểm tra dữ liệu thiếu, \texttt{groupby}, \texttt{query}, tương quan & \texttt{pipeline.py}, notebook \\
scikit-learn & \texttt{Pipeline}, \texttt{ColumnTransformer}, \texttt{SimpleImputer}, \texttt{StandardScaler}, \texttt{OneHotEncoder}, \texttt{Ridge}, \texttt{RandomForestRegressor}, \texttt{DummyRegressor}, chia tập và các độ đo & \texttt{train\_model.py} \\
joblib & Lưu mô hình đã huấn luyện & \texttt{train\_model.py} \\
Matplotlib, Seaborn & Biểu đồ & \texttt{pipeline.py}, \texttt{train\_model.py} \\
Jupyter & Notebook trình bày các bước phân tích & \texttt{melbourne\_\allowbreak data\_\allowbreak analysis\_\allowbreak pipeline.ipynb} \\
\bottomrule
\end{tabular}
\nguon{\texttt{requirements.txt}; việc dùng thư viện trong từng tệp xem mã nguồn.}
\end{table}
